\documentclass[12pt, a4paper]{article}

\usepackage[utf8]{inputenc}
\usepackage[T1]{fontenc}
\usepackage{amsmath, amsthm, amssymb}
\usepackage{mathtools}
\usepackage{geometry}
\usepackage{hyperref}
\usepackage{enumitem}
\usepackage{xcolor}
\usepackage{mdframed}

\geometry{margin=2.5cm}

\theoremstyle{definition}
\newtheorem{definition}{Definition}[section]
\newtheorem{example}[definition]{Example}

\theoremstyle{plain}
\newtheorem{theorem}[definition]{Theorem}
\newtheorem{lemma}[definition]{Lemma}
\newtheorem{corollary}[definition]{Corollary}

\theoremstyle{remark}
\newtheorem{remark}[definition]{Remark}

\newmdenv[
  backgroundcolor=blue!5, linecolor=blue!30, linewidth=1pt,
  innertopmargin=8pt, innerbottommargin=8pt, skipabove=10pt, skipbelow=6pt
]{bluebox}
\newmdenv[
  backgroundcolor=orange!5, linecolor=orange!40, linewidth=1pt,
  innertopmargin=8pt, innerbottommargin=8pt, skipabove=10pt, skipbelow=6pt
]{orangebox}
\newmdenv[
  backgroundcolor=green!5, linecolor=green!40, linewidth=1pt,
  innertopmargin=8pt, innerbottommargin=8pt, skipabove=10pt, skipbelow=6pt
]{greenbox}

\newcommand{\CM}{C_M}
\newcommand{\CW}{C_W}
\newcommand{\SM}{S^M}
\newcommand{\SW}{S^W}
\newcommand{\school}[2]{\mathrm{sch}(#1,#2)}
\newcommand{\stud}[2]{\mathrm{st}(#1,#2)}
\newcommand{\topK}{\mathrm{topK}}
\newcommand{\topOne}{\mathrm{topOne}}
\newcommand{\quota}{b}

\title{\textbf{Classical and Fleiner Stability:\\
Equivalence of Definitions}\\[0.5em]
\large Complete Mathematical Proof}
\author{}
\date{}

\begin{document}
\maketitle
\tableofcontents
\newpage

%=============================================================================
\section{Model and Definitions}
%=============================================================================

\subsection{Problem data}

Fix:
\begin{itemize}
  \item finite sets of schools $M$ and students $W$;
  \item a set of admissible pairs $E \subseteq M \times W$ (the bipartite graph);
  \item for each school $m \in M$: a strict linear order $\succ_m$ on $W$
    and a quota $\quota(m) \in \mathbb{N}$;
  \item for each student $w \in W$: a strict linear order $\succ_w$ on $M$.
\end{itemize}

The notation $a \succ b$ means ``$a$ is preferred to $b$''.

\subsection{Notation}

For $S \subseteq E$ write:
\begin{align*}
  \school{S}{m} &:= \{w \in W : (m,w) \in S\}
    \quad\text{--- students assigned to school } m \text{ in } S,\\
  \stud{S}{w}   &:= \{m \in M : (m,w) \in S\}
    \quad\text{--- school(s) of student } w \text{ in } S.
\end{align*}
Since $S$ is a matching, $|\stud{S}{w}| \leq 1$.

The choice functions are defined via ranking:
\begin{align*}
  \CM(A) &:= \bigl\{(m,w) \in A : w \in \topK(\succ_m,\,\quota(m),\,\school{A}{m})\bigr\},\\
  \CW(A) &:= \bigl\{(m,w) \in A : m \in \topOne(\succ_w,\,\stud{A}{w})\bigr\},
\end{align*}
where $\topK(\succ, k, B)$ denotes the $k$ best elements of $B$ under $\succ$,
and $\topOne(\succ, B)$ denotes the unique best element of $B$.

\subsection{Two definitions of stability}

\begin{definition}[Classical stability]
\label{def:classical}
A matching $S \subseteq E$ is \emph{classically stable} if it is individually
rational (i.e.\ $|\school{S}{m}| \leq \quota(m)$ and $|\stud{S}{w}| \leq 1$
for all $m, w$) and has no blocking pair.

A \emph{blocking pair} is a pair $(m, w) \in E \setminus S$ such that
simultaneously:
\begin{itemize}
  \item \textbf{school $m$ wants $w$:} $|\school{S}{m}| < \quota(m)$, or
    $\exists\, w' \in \school{S}{m} : w \succ_m w'$;
  \item \textbf{student $w$ wants $m$:} $\stud{S}{w} = \emptyset$, or
    $\exists\, m' \in \stud{S}{w} : m \succ_w m'$.
\end{itemize}
\end{definition}

\begin{definition}[Fleiner stability]
\label{def:fleiner}
A pair $(\SM, \SW)$ of subsets of $E$ is a \emph{stable pair} if
\begin{align}
  \SM \cup \SW &= E, \tag{F1}\label{F1}\\
  \CM(\SM)     &= \SM \cap \SW, \tag{F2}\label{F2}\\
  \CW(\SW)     &= \SM \cap \SW. \tag{F3}\label{F3}
\end{align}
A matching $S \subseteq E$ is \emph{Fleiner-stable} if $S = \SM \cap \SW$
for some stable pair $(\SM, \SW)$.
\end{definition}

\begin{bluebox}
\textbf{Intuition.}
$\SM$ represents ``proposals from schools'' (pairs that schools are willing to
consider). $\SW$ represents ``proposals from students''.
Equation~\eqref{F2} says: what schools actually \emph{choose} from their
proposals $\SM$ equals the final matching $S = \SM \cap \SW$.
Equation~\eqref{F3} is the symmetric statement for students.
\end{bluebox}

\subsection{Main theorem}

\begin{theorem}[Equivalence of definitions]
\label{thm:equiv}
A matching $S \subseteq E$ is classically stable if and only if it is
Fleiner-stable.
\end{theorem}

The proof is split into two directions. The direction $(\Leftarrow)$ (Fleiner
$\Rightarrow$ classical) is easier and shorter; the direction $(\Rightarrow)$
(classical $\Rightarrow$ Fleiner) requires an explicit construction of a stable pair.

%=============================================================================
\section{Direction $(\Leftarrow)$: Fleiner $\Rightarrow$ Classical}
%=============================================================================

\textbf{Given:} $S = \SM \cap \SW$ where $(\SM, \SW)$ is a stable pair.

\textbf{To show:} $S$ is classically stable, i.e.\ it is a matching and has
no blocking pair.

\subsection{Individual rationality}

\begin{lemma}
\label{lem:ir}
$S = \SM \cap \SW$ is a matching.
\end{lemma}
\begin{proof}
By~\eqref{F2}: $S = \CM(\SM)$. By definition $\CM$ includes for each school
$m$ exactly $\topK(\succ_m, \quota(m), \school{\SM}{m})$ students, hence at
most $\quota(m)$. Therefore $|\school{S}{m}| \leq \quota(m)$.

By~\eqref{F3}: $S = \CW(\SW)$. By definition $\CW$ includes for each student
$w$ exactly $\topOne(\succ_w, \stud{\SW}{w})$, hence at most one school.
Therefore $|\stud{S}{w}| \leq 1$.
\end{proof}

\subsection{Absence of blocking pairs}

\begin{lemma}
\label{lem:no_block}
$S$ has no blocking pair.
\end{lemma}
\begin{proof}
Suppose for contradiction that $(m, w) \in E \setminus S$ is a blocking pair.
By~\eqref{F1}: $(m, w) \in \SM \cup \SW$, so $(m, w) \in \SM$ or
$(m, w) \in \SW$.

\medskip
\noindent\textbf{Claim A: $(m, w) \in \SM$.}

Suppose $(m, w) \notin \SM$. Then $(m, w) \in \SW \setminus \SM$.
Since $(m, w) \notin S = \CW(\SW)$, we have
$m \notin \topOne(\succ_w, \stud{\SW}{w})$.
Yet $m \in \stud{\SW}{w}$ (since $(m, w) \in \SW$), so $\stud{\SW}{w}$
is non-empty and has a unique best element $m^* \in \stud{\SW}{w}$; since $m$
is not the best, we get $m^* \succ_w m$.

Since $m^* \in \topOne(\succ_w, \stud{\SW}{w})$, we have
$(m^*, w) \in \CW(\SW) = S$.
Therefore $m^* \in \stud{S}{w}$, i.e.\ $w$ is already assigned to $m^*$.

Now consider ``student $w$ wants $m$'':
\begin{itemize}
  \item If $\stud{S}{w} = \emptyset$: contradiction with $m^* \in \stud{S}{w}$.
  \item If $\exists\, m' \in \stud{S}{w} : m \succ_w m'$: since
    $\stud{S}{w} \ni m^*$, the only option is $m' = m^*$, giving
    $m \succ_w m^*$. But $m^* \succ_w m$ --- contradiction (transitivity
    then gives $m \succ_w m$, violating irreflexivity).
\end{itemize}
Both cases yield a contradiction. Hence $(m, w) \in \SM$.

\medskip
\noindent\textbf{Claim B: $(m, w) \in \SW$.}

Suppose $(m, w) \notin \SW$. Then $(m, w) \in \SM \setminus \SW$.
Since $(m, w) \notin S = \CM(\SM)$, we have
$w \notin \topK(\succ_m, \quota(m), \school{\SM}{m})$.
Yet $w \in \school{\SM}{m}$ (since $(m, w) \in \SM$).

By definition of $\topK$: since $w \in \school{\SM}{m}$ and $w \notin \topK$,
there are at least $\quota(m)$ students in $\school{\SM}{m}$ preferred to $w$:
\[
  \bigl|\{w' \in \school{\SM}{m} : w' \succ_m w\}\bigr| \geq \quota(m).
\]
Consequently $|\topK(\succ_m, \quota(m), \school{\SM}{m})| = \quota(m)$,
i.e.\ $|\school{S}{m}| = \quota(m)$ (the quota is full), and every student
in $\school{S}{m} = \topK$ is preferred to $w$:
$\forall w'' \in \school{S}{m} : w'' \succ_m w$.

Now consider ``school $m$ wants $w$'':
\begin{itemize}
  \item If $|\school{S}{m}| < \quota(m)$: contradiction with
    $|\school{S}{m}| = \quota(m)$.
  \item If $\exists\, w' \in \school{S}{m} : w \succ_m w'$: but we showed
    $w'' \succ_m w$ for all $w'' \in \school{S}{m}$, in particular for $w'$.
    So $w \succ_m w'$ and $w' \succ_m w$ --- contradiction with irreflexivity.
\end{itemize}
Both cases yield a contradiction. Hence $(m, w) \in \SW$.

\medskip
\noindent\textbf{Conclusion.}
By both claims $(m, w) \in \SM \cap \SW = S$, contradicting $(m, w) \notin S$.
\end{proof}

\begin{corollary}
Every Fleiner-stable matching is classically stable.
\end{corollary}

\begin{orangebox}
\textbf{Key idea of direction $(\Leftarrow)$.}
The Fleiner equations algebraically encode why blocking pairs are impossible:
if $(m, w) \notin S$, then either $w$ is displaced in $\SM$ (the quota of $m$
is filled by better students --- Claim~B), or $m$ is displaced in $\SW$
(student $w$ is already assigned to a better school --- Claim~A). Both claims
together give $(m, w) \in \SM \cap \SW = S$, which is impossible.
\end{orangebox}

%=============================================================================
\section{Direction $(\Rightarrow)$: Classical $\Rightarrow$ Fleiner}
%=============================================================================

\textbf{Given:} $S \subseteq E$ is classically stable.

\textbf{To show:} construct $(\SM, \SW)$ explicitly such that $S = \SM \cap \SW$
and conditions~\eqref{F1}--\eqref{F3} hold.

\subsection{Two types of ``non-membership'' in a matching}

For each pair $(m, w) \in E \setminus S$ (a pair that did not enter the
matching) we introduce two types:

\begin{definition}
$(m, w) \in E \setminus S$ is \emph{school-rejected} if
\[
  |\school{S}{m}| = \quota(m) \quad\text{and}\quad
  \forall\, w' \in \school{S}{m} : w' \succ_m w.
\]
In other words: school $m$'s quota is full, and every admitted student is
preferred to $w$.
\end{definition}

\begin{definition}
$(m, w) \in E \setminus S$ is \emph{student-rejected} if
\[
  \stud{S}{w} \neq \emptyset \quad\text{and}\quad
  \forall\, m' \in \stud{S}{w} : m' \succ_w m.
\]
In other words: student $w$ is already assigned to a school that she prefers
to $m$.
\end{definition}

\begin{lemma}[Coverage]
\label{lem:coverage}
Every pair $(m, w) \in E \setminus S$ is school-rejected or student-rejected.
\end{lemma}
\begin{proof}
Since $S$ is classically stable, $(m, w)$ is not a blocking pair. Therefore
at least one of the two blocking conditions fails:
\begin{itemize}
  \item ``School $m$ does not want $w$'': $|\school{S}{m}| = \quota(m)$ and
    $\forall\, w' \in \school{S}{m} : w' \succ_m w$ --- exactly the definition
    of school-rejected.
  \item ``Student $w$ does not want $m$'': $\stud{S}{w} \neq \emptyset$ and
    $\forall\, m' \in \stud{S}{w} : m' \succ_w m$ --- exactly the definition
    of student-rejected.\qedhere
\end{itemize}
\end{proof}

\subsection{Construction of the stable pair}

\begin{definition}[Canonical stable pair]
\label{def:canonical}
Define:
\begin{align}
  \SW &:= S \;\cup\; \{(m,w) \in E \setminus S : (m,w)\text{ is student-rejected}\},
  \label{eq:SW}\\
  \SM &:= S \;\cup\; \{(m,w) \in E \setminus S : (m,w)\text{ is school-rejected
    but not student-rejected}\}.
  \label{eq:SM}
\end{align}
\end{definition}

\begin{bluebox}
\textbf{Partition idea.}
The pairs in $E \setminus S$ fall into three groups:
\begin{enumerate}
  \item Only school-rejected (quota full, but $w$ is unmatched or not yet
    placed better): $\to \SM$.
  \item Only student-rejected (student placed better, but $m$'s quota is not
    full of better students): $\to \SW$.
  \item Both types simultaneously (quota full with better students AND student
    placed better): $\to \SW$.
\end{enumerate}
Thus every pair falls into at least one set, and only pairs of type~1 end up
exclusively in $\SM$.
\end{bluebox}

\subsection{Verification of condition~\eqref{F1}: $\SM \cup \SW = E$}

\begin{lemma}
$\SM \cup \SW = E$.
\end{lemma}
\begin{proof}
The inclusion $\SM \cup \SW \subseteq E$ is immediate (all pairs are taken
from $E$).

For the reverse: take $(m, w) \in E$.
\begin{itemize}
  \item If $(m, w) \in S$: then $(m, w) \in S \subseteq \SM$.
  \item If $(m, w) \in E \setminus S$: by Lemma~\ref{lem:coverage} the pair
    is school-rejected or student-rejected. By construction:
    if student-rejected then $(m, w) \in \SW$;
    if only school-rejected then $(m, w) \in \SM$.\qedhere
\end{itemize}
\end{proof}

\subsection{Verification of the intersection: $\SM \cap \SW = S$}

\begin{lemma}
\label{lem:intersection}
$\SM \cap \SW = S$.
\end{lemma}
\begin{proof}
$S \subseteq \SM \cap \SW$ by construction (both sets contain $S$).

We show $\SM \cap \SW \subseteq S$. Let $(m, w) \in \SM \cap \SW$ with
$(m, w) \notin S$.
\begin{itemize}
  \item $(m, w) \in \SM \setminus S$: the pair is school-rejected and
    \emph{not} student-rejected. In particular: $\stud{S}{w} = \emptyset$ or
    $\exists\, m' \in \stud{S}{w} : m \succ_w m'$.
  \item $(m, w) \in \SW \setminus S$: the pair is student-rejected, i.e.\
    $\stud{S}{w} \neq \emptyset$ and $\forall\, m' \in \stud{S}{w} : m' \succ_w m$.
\end{itemize}
From the first bullet, $(m, w)$ is not student-rejected; from the second
bullet, it is student-rejected. Contradiction. Hence $(m, w) \in S$.
\end{proof}

\subsection{Verification of condition~\eqref{F2}: $\CM(\SM) = S$}

\begin{lemma}
\label{lem:choiceM}
$\CM(\SM) = S$.
\end{lemma}
\begin{proof}
Fix a school $m$. We need to show:
\[
  \topK\bigl(\succ_m,\,\quota(m),\,\school{\SM}{m}\bigr) = \school{S}{m}.
\]

\textbf{Step 1: Describe $\school{\SM}{m}$.}

$\school{\SM}{m} = \school{S}{m} \cup \{w' : (m, w') \in \SM \setminus S\}$.

By construction~\eqref{eq:SM}: $(m, w') \in \SM \setminus S$ means the pair
$(m, w')$ is school-rejected. By definition of school-rejected:
$|\school{S}{m}| = \quota(m)$ and $\forall\, w'' \in \school{S}{m} : w'' \succ_m w'$.
In other words, every student in $\school{S}{m}$ is preferred to $w'$.

Therefore: every element of $\school{\SM}{m} \setminus \school{S}{m}$ is
strictly worse than every element of $\school{S}{m}$.

\textbf{Step 2: Case $|\school{S}{m}| = \quota(m)$.}

Then $\school{\SM}{m} = \school{S}{m} \cup \{\text{students worse than all
of } \school{S}{m}\}$. The best $\quota(m)$ students in $\school{\SM}{m}$
are exactly $\school{S}{m}$:
\[
  \topK\bigl(\succ_m, \quota(m), \school{\SM}{m}\bigr) = \school{S}{m}.
  \quad\checkmark
\]

\textbf{Step 3: Case $|\school{S}{m}| < \quota(m)$.}

The school-rejected condition requires $|\school{S}{m}| = \quota(m)$, which
fails here, so no pair is added to $\SM$ for school $m$.
Hence $\school{\SM}{m} = \school{S}{m}$. Since $|\school{S}{m}| < \quota(m)$:
\[
  \topK\bigl(\succ_m, \quota(m), \school{S}{m}\bigr) = \school{S}{m}.
  \quad\checkmark
\]

In both cases $\topK(\succ_m, \quota(m), \school{\SM}{m}) = \school{S}{m}$,
which means
$\CM(\SM) \cap \{(m, w) : w \in W\} = S \cap \{(m, w) : w \in W\}$
for every $m$. Summing over all $m$: $\CM(\SM) = S$.
\end{proof}

\subsection{Verification of condition~\eqref{F3}: $\CW(\SW) = S$}

\begin{lemma}
\label{lem:choiceW}
$\CW(\SW) = S$.
\end{lemma}
\begin{proof}
Fix a student $w$. We need to show:
\[
  \topOne\bigl(\succ_w,\,\stud{\SW}{w}\bigr) = \stud{S}{w}.
\]

\textbf{Case 1: $\stud{S}{w} = \emptyset$ (student is unmatched).}

By construction~\eqref{eq:SW}: $(m', w) \in \SW \setminus S$ requires
$(m', w)$ to be student-rejected, i.e.\ $\stud{S}{w} \neq \emptyset$.
Since $\stud{S}{w} = \emptyset$, no pair $(m', w)$ enters $\SW \setminus S$.
Hence $\stud{\SW}{w} = \emptyset$ and
$\topOne(\succ_w, \emptyset) = \emptyset = \stud{S}{w}$. $\checkmark$

\textbf{Case 2: $\stud{S}{w} = \{m^*\}$ (student is matched to school $m^*$).}

Describe $\stud{\SW}{w}$:
\[
  \stud{\SW}{w} = \{m^*\} \cup
  \bigl\{m' : (m', w) \in E \setminus S,\;
  (m', w)\text{ is student-rejected}\bigr\}.
\]
For every student-rejected pair $(m', w)$ the definition gives
$\forall\, m'' \in \stud{S}{w} : m'' \succ_w m'$, i.e.\ $m^* \succ_w m'$
(since $\stud{S}{w} = \{m^*\}$). Hence:
\[
  \stud{\SW}{w}
  = \{m^*\} \cup \{m' : m^* \succ_w m'\text{ for some pair in }E \setminus S\}
  \subseteq \{m \in M : m^* \succ_w m \text{ or } m = m^*\}.
\]
In other words, $m^*$ is at least as good (under $\succ_w$) as every element
of $\stud{\SW}{w}$. Therefore $m^* \in \topOne(\succ_w, \stud{\SW}{w})$, and
since $\topOne$ returns exactly one element,
$\topOne(\succ_w, \stud{\SW}{w}) = \{m^*\} = \stud{S}{w}$. $\checkmark$

\medskip
In both cases $\topOne(\succ_w, \stud{\SW}{w}) = \stud{S}{w}$ for every $w$,
giving $\CW(\SW) = S$.
\end{proof}

\subsection{Conclusion of direction $(\Rightarrow)$}

\begin{theorem}
\label{thm:classical_to_fleiner}
Every classically stable matching is Fleiner-stable.
\end{theorem}
\begin{proof}
For a given classically stable $S$, construct $(\SM, \SW)$ as in
Definition~\ref{def:canonical}. By the lemmas above:
\begin{itemize}
  \item $\SM \cup \SW = E$ (\ref{F1}): Lemma~\ref{lem:coverage} + construction;
  \item $\CM(\SM) = S$ and $\CW(\SW) = S$: Lemmas~\ref{lem:choiceM}
    and~\ref{lem:choiceW};
  \item $\SM \cap \SW = S$: Lemma~\ref{lem:intersection}.
\end{itemize}
Therefore $\CM(\SM) = S = \SM \cap \SW$ and $\CW(\SW) = S = \SM \cap \SW$,
so conditions~\eqref{F2} and~\eqref{F3} hold. Thus $(\SM, \SW)$ is a stable
pair and $S = \SM \cap \SW$ is Fleiner-stable.
\end{proof}

%=============================================================================
\section{Main Theorem and Corollaries}
%=============================================================================

\begin{theorem}[Equivalence, full statement]
\label{thm:main_full}
Let $M$, $W$ be finite sets, $E \subseteq M \times W$, with a strict linear
preference order for each agent and quotas $\quota(m)$ for schools.
Then the following conditions on $S \subseteq E$ are equivalent:
\begin{enumerate}[label=(\alph*)]
  \item $S$ is \textbf{classically stable}: a matching with no blocking pair;
  \item $S$ is \textbf{Fleiner-stable}: there exists $(\SM, \SW)$ with
    $\SM \cup \SW = E$, $\CM(\SM) = \SM \cap \SW$, $\CW(\SW) = \SM \cap \SW$,
    and $S = \SM \cap \SW$.
\end{enumerate}
\end{theorem}
\begin{proof}
$(b) \Rightarrow (a)$: Lemmas~\ref{lem:ir} and~\ref{lem:no_block}.\\
$(a) \Rightarrow (b)$: Theorem~\ref{thm:classical_to_fleiner}.
\end{proof}

\begin{greenbox}
\textbf{Practical consequence.}
All theorems proved in the Fleiner language --- existence of a stable matching,
lattice structure, the Rural Hospital Theorem --- transfer automatically to
classically stable matchings via this equivalence.
\end{greenbox}

\begin{corollary}[Existence of a classically stable matching]
For any finite $M$, $W$, linear preferences, and quotas, a classically stable
matching exists.
\end{corollary}
\begin{proof}
By the existence theorem (in the Fleiner language) a Fleiner-stable $S$ exists.
By Theorem~\ref{thm:main_full} it is classically stable.
\end{proof}

\begin{corollary}[Rural Hospital Theorem, classical formulation]
If $S_1$ and $S_2$ are classically stable and $|\school{S_1}{m}| < \quota(m)$,
then $\school{S_1}{m} = \school{S_2}{m}$.
\end{corollary}
\begin{proof}
Translate to the Fleiner language via Theorem~\ref{thm:main_full}, apply the
Rural Hospital Theorem, and translate back.
\end{proof}

%=============================================================================
\section{Asymmetry of the Two Directions}
%=============================================================================

The two directions of the proof differ fundamentally in nature.

\textbf{Direction $(\Leftarrow)$} (Fleiner $\to$ classical) is a
\emph{verification}: the given object (a stable pair) satisfies a condition
(no blocking pairs). The proof is purely logical: the equations~\eqref{F2}
and~\eqref{F3} directly explain why every potential blocking pair is impossible.
The two cases (school side and student side) correspond exactly to the two
equations.

\textbf{Direction $(\Rightarrow)$} (classical $\to$ Fleiner) is a
\emph{construction}: one must \emph{build} an object (a stable pair) and then
verify three equations. The construction is natural: $\SW$ contains
``student-rejected'' pairs, $\SM$ contains ``school-rejected but not
student-rejected'' pairs. The key observation making the construction possible
is that every pair in $E \setminus S$ is rejected by at least one side
(otherwise it would be a blocking pair). The key observation for
$\SM \cap \SW = S$ is that membership in $\SM \setminus S$ and membership in
$\SW \setminus S$ are mutually exclusive conditions: a pair in $\SM \setminus S$
is by definition not student-rejected, while a pair in $\SW \setminus S$ is by
definition student-rejected.

\textbf{Status in the formalization.}
In the file \texttt{ClassicalBridge.lean}, direction $(\Leftarrow)$ is fully
proved in Lean~4 (no \texttt{sorry}). Direction $(\Rightarrow)$ is marked
\texttt{sorry}: the construction above is mathematically correct (as shown in
this document), but its formal verification has not yet been completed.

\end{document}